\section{Experimental assessment}
\label{sec:tests}

The assessment comprises three campaigns based on published geysering
experiments. Simulations start from the
documented experimental initial conditions, except for the explicitly
identified C9 phase-one all-water variant used in the analytic comparison.
Shaft motion, gas release, and regime selection are outcomes of the production
calculations. Outside the explicitly labelled Test~A sensitivity and the
disclosed Test~B screen, the calculations use no prescribed release rate,
scripted geyser height, overflow accumulator, or per-case observation fit.

The three campaigns provide complementary evidence.
Campaign~1 (Section~\ref{sec:tests:vw2011}) reproduces the water-filled
ventilation-shaft experiments of Vasconcelos and Wright
\cite{vasconcelos2011}, a widely used controlled benchmark for geysering
models. It tests time-resolved histories in both regimes: Test~A represents the
non-geysering branch, and Test~B represents the geysering branch.
Campaign~2 (Section~\ref{sec:tests:cong2017}) reproduces the
horizontal-pipe/vertical-riser experiments of Cong, Chan and Lee
\cite{cong2017}. Because that study condenses its outcomes into an explicit
two-parameter geysering criterion, this campaign tests the model as a regime
classifier across a systematic parameter sweep rather than against individual
histories.
Campaign~3 (Section~\ref{sec:tests:liu2020}) addresses the
prototype-motivated junction-chamber experiments of Liu, Shao and Zhu
\cite{liu2020}. These experiments connect the laboratory benchmarks to a
field-scale dropshaft geometry.
Other studies are used as corroborating references rather than reproduction
targets. They include the field observations of \cite{wright2011note}, the
larger-scale releases designed for three-dimensional CFD calibration
\cite{muller2017}, the covered-manhole geysers of \cite{zhang2022}, the
mitigation study of \cite{qian2020}, and the vertical-pipe violent geysers of
\cite{leon2018}. Their geometries are either insufficiently controlled for a
one-dimensional comparison, or their focus (cover dynamics, mitigation
devices, or flashing-driven eruption) lies outside the entrapped-air release
mechanism considered here.

\subsection{Campaign 1: water-filled ventilation shafts
(Vasconcelos and Wright, 2011)}
\label{sec:tests:vw2011}

Vasconcelos and Wright \cite{vasconcelos2011} used a horizontal acrylic pipe
of diameter $D=0.094$~m with both far ends closed. A ventilation tower of
length $L=0.610$~m and variable diameter $D_t$ was mounted on a coupling
$3.516$~m from the upstream end. A butterfly valve separated a
$0.546$-m-long upstream air chamber from the water-filled reach
(Fig.~\ref{fig:apparatus}). The air chamber was pressurized to an initial gauge
head $H_{a0}$, and the tower was partially filled to an initial level
$Y_{fs,0}$. Each run began when the valve was opened by hand in less than one
second. The pressurized pocket then propagated along the pipe crown toward the
tower and vented through the water column. A pressure transducer was located
$1.07$~m downstream of the valve. The experimental programme varied
$D_t\in\{57.1,\,44.4,\,25.4,\,12.7\}$~mm ($D_t/D=0.607$--$0.135$),
$H_{a0}\in\{0.305,\,0.610,\,0.915\}$~m, and
$Y_{fs,0}\in\{0.254,\,0.356,\,0.457\}$~m. Tower diameter controlled the
observed response: for $D_t/D=0.607$, the free-surface rise never exceeded
10\% of $L$, whereas for $D_t/D=0.135$, the free surface reached the tower top
in every repetition.

Two configurations are selected to span the two regimes
(Table~\ref{tab:cases}). Following \cite{vasconcelos2011}, results are
non-dimensionalized by the tower length and the tower gravity-velocity
scale: elevations and pressure heads as $Y^{*}=Y/L$ and $H^{*}=H/L$, and time
as $T_{\mathrm{rel}}^{*}=t\sqrt{gD_t}/L$. Geysering is defined as the free surface
$Y^{*}_{fs}$ reaching the tower top ($Y^{*}=1$) before the rising air--water
interface $Y^{*}_{int}$ catches the free surface.

\begin{table}[htbp]
\centering
\caption{Selected test configurations from the experimental matrix of
\cite{vasconcelos2011}.}
\label{tab:cases}
\begingroup
\small
\setlength{\tabcolsep}{3.5pt}
\begin{tabular}{@{}lcccccc@{}}
\toprule
Test & \shortstack{$D_t$\\(mm)} & $D_t/D$ & \shortstack{$H_{a0}$\\(m)} &
\shortstack{$Y_{fs,0}$\\(m)} & \shortstack{$L/\sqrt{gD_t}$\\(s)} &
\shortstack{Observed\\regime} \\
\midrule
A & 57.1 & 0.607 & 0.305 & 0.356 & 0.815 & no geysering \\
B & 12.7 & 0.135 & 0.610 & 0.356 & 1.728 & geysering (3/3) \\
\bottomrule
\end{tabular}
\endgroup
\end{table}

The one-dimensional model geometry follows the apparatus exactly: upstream
air chamber $0.546$~m (the initial pocket volume,
$3.79\times10^{-3}$~m$^3$), middle pipe $2.970$~m, downstream closed pipe
$0.490$~m, tower at $x=3.516$~m, and the transducer sampled at $x=1.616$~m.
The initial pocket pressure is
$P_{a0}=P_{\mathrm{atm}}+\rho_l g H_{a0}$; the pipe downstream of the valve
and the tower below $Y_{fs,0}$ are initially at rest and water-filled. The
experimental curves used for comparison were digitized from the published
figures of \cite{vasconcelos2011} (pressure head from Figs.~5 and~6;
free-surface and interface trajectories from Figs.~7 and~8, three
experimental repetitions each). Pressure repetitions are summarized by a
representative pointwise curve, while level repetitions are redrawn
individually as vector markers.

\subsubsection{Test A: wide tower, non-geysering branch
($D_t/D=0.607$)}
\label{sec:tests:vw2011:caseA}

Test~A is the wide-tower, non-geysering member of Campaign~1. It first tests
whether the reduced model selects venting without overflow. It then compares
the computed pressure and tower-level scales with the measurements. In the
experiment, the trapped pocket vents through the tower before the free surface
reaches the rim. During interface rise, $Y^{*}_{fs}$ remains close to 0.6, and
the transducer head retains a long plateau until pocket exhaustion
\citep{vasconcelos2011}. The archived one-dimensional and planar OpenFOAM
calculations select the same branch. Two complementary archives support the
comparison. Figure~\ref{fig:caseA_snapshots} examines common-clock event
ordering and phase topology. Figure~\ref{fig:caseA_curves} quantifies a
separately declared connected-pocket closure sensitivity. The two figures are
not presented as outputs from one numerical realization.

The nearest-time states in Fig.~\ref{fig:caseA_snapshots} divide the event into
four stages. Near $t=1.50$~s, both calculations show the early post-valve
transient while the pocket remains in the horizontal pipe. Near $t=6.90$~s,
gas has reached the junction and entered the water-filled tower. Near
$t=7.60$~s, the breakthrough transition is accompanied by rapid drainage and
phase redistribution in the tower. The one-dimensional liquid-equivalent
height has fallen to $0.150$~m, while the planar field resolves the irregular
junction-scale breakup that a cross-section-averaged column cannot represent.
By $t=10.10$~s, the one-dimensional archive retains only $0.0155$~m of
liquid-equivalent height, concentrated visibly below approximately $0.082$~m
above the tower base; the planar field retains a broader, laterally deformed
residual-liquid region. The paired 1D and 2D archives are sampled without a
time shift, with a maximum time mismatch of $0.020$~s. They preserve
horizontal-pocket propagation, tower entry, and breakthrough ordering and
remain on the non-geysering branch, but the late phase distribution is not
matched. The snapshots therefore support branch selection and event ordering,
not a quantitative phase-field error norm. For legibility only, 1D tower cells
with $\alpha_l\leq0.02$ are omitted from the raster rendering; their liquid
volume remains included in the reported equivalent height.

\begin{figure}[htbp]
\centering
\includegraphics[width=\textwidth]{caseA_1d2d_snapshots_4frame.pdf}
\caption{Test~A phase evolution near $t=1.50$, $6.90$, $7.60$, and
$10.10$~s: present 1D model (left) and planar OpenFOAM simulation (right).
The corresponding archived 1D times are $1.480$, $6.908$, $7.608$, and
$10.108$~s, respectively; no time shift is applied. Blue denotes water and
white denotes gas. Row~(c) shows the breakthrough transition and rapid tower
drainage, and row~(d) shows the small residual-liquid region near the tower
base. In the 1D rendering only, cells with $\alpha_l\leq0.02$ are visually
omitted while their liquid volume remains in the equivalent-height audit.}
\label{fig:caseA_snapshots}
\end{figure}

The native-time histories in Fig.~\ref{fig:caseA_curves} provide the
quantitative comparison. The one-dimensional history is an explicitly
isolated sensitivity. It combines the contact-preserving release update with a
connected-pocket pressure-relaxation closure of time scale $\tau_b=0.10$~s.
This optional closure probes unresolved exchange between the coherent pocket
and the draining wall film, and it is not used for the other campaign claims.
The values below characterize the disclosed sensitivity rather than an
independently validated production closure.

The pre-venting pressure scale is close to the experiment, but the terminal
release remains late. Over
$4\leq T_{\mathrm{rel}}^{*}\leq7.5$, the one-dimensional plateau median is
$H^{*}=0.539$, compared with $0.537$ for the extracted experimental central
trace. After hydrostatic conversion from the probe elevation, the
crown-referenced two-dimensional mean is $0.562$ over
$1\leq T_{\mathrm{rel}}^{*}\leq7$. This datum conversion is not a fitted
offset. The main discrepancy lies in the release clock. The experimental,
two-dimensional, and one-dimensional traces first fall below $H^{*}=0.30$ at
$T_{\mathrm{rel}}^{*}=8.28$, 8.63, and 9.20, respectively. The
one-dimensional sensitivity therefore preserves the plateau but delays
terminal depressurization by about $0.92$ dimensionless time units. No time or
pressure shift is applied.

The short one-dimensional spike near $T_{\mathrm{rel}}^{*}=9.18$ coincides
with the archived level-coalescence event. It is retained for provenance but
is not interpreted as a resolved physical pressure rebound or included in the
plateau statistic.

All three records remain below the rim, and the one-dimensional level-event
times fall within the experimental spread.
The one-dimensional free surface rises gradually to
$Y^{*}_{fs,\max}=0.613$. The largest digitized
experimental marker is $0.649$, and the two-dimensional maximum is $0.623$;
none reaches the rim. The one-dimensional interface crosses
$Y^{*}_{int}=0.05$ at $T_{\mathrm{rel}}^{*}=7.85$. The raw archive reaches
exact equality with the free surface at 9.18, whereas the displayed trace stops
at the one-cell criterion at 9.13. Across the three digitized repetitions, the
first markers above the same $Y^{*}_{int}=0.05$ threshold span approximately
7.36--7.94. Piecewise-linear crossings of the digitized interface and
free-surface tracks span approximately 8.85--9.19. Both one-dimensional event
times lie within the observed run-to-run spread, although the sensitivity does
not reproduce the fastest interface rise. The sensitivity captures the
non-geysering outcome, pressure plateau, and maximum free-surface scale, but
terminal depressurization remains late. The remaining shape differences are
consistent with unresolved coherent-core/wall-film exchange. These data do not
identify that mechanism uniquely.

\begin{figure}[htbp]
\centering
\includegraphics[width=\textwidth]{caseA_experiment_1d2d_curves.pdf}
\caption{Test~A comparison with \cite{vasconcelos2011}: (a) pressure head and
(b) normalized free-surface and interface elevations. Black denotes the
experiment, orange the present 1D model, and blue OpenFOAM.}
\label{fig:caseA_curves}
\end{figure}

The planar OpenFOAM calculation is a supporting spatial diagnostic rather than
a geometry-exact reference. In the vertical plane, the tower-to-pipe area
ratio is $D_t/D=0.607$. The circular apparatus instead has
$A_t/A_p=(D_t/D)^2=0.369$. The planar model also cannot reproduce a
three-dimensional annular wall film. Quantitative assessment is therefore
anchored to the experiment.
Within that scope, Test~A supports the wide-tower branch and the bulk pressure
and water-level scales, while revealing a systematic late bias in terminal
pocket depressurization. The interface event times remain within the digitized
three-run spread. This non-geysering reference provides the controlled contrast
for the narrow-tower Test~B; it is not an independent validation of the
optional connected-pocket closure.

\FloatBarrier

\subsubsection{Test B: small tower, geysering branch
($D_t/D=0.135$)}
\label{sec:tests:vw2011:caseB}

The narrow tower provides the geysering counterpart to Test~A. Vasconcelos and
Wright reported a geyser in every repetition. In the three Fig.~8 records, gas
first appears in the tower over $T_{\mathrm{rel}}^{*}=3.648$--$3.742$. The free
surface reaches the rim over $T_{\mathrm{rel}}^{*}=3.900$--$3.981$ and therefore
precedes gas breakthrough. The credible high-interface markers occur over
$T_{\mathrm{rel}}^{*}=4.09$--$4.17$. The Fig.~6 transducer record has a
pre-release plateau near $H^{*}=0.76$ and a collapse envelope of approximately
$T_{\mathrm{rel}}^{*}=4.02$--$4.18$ \cite{vasconcelos2011}.
The $D_t/D=0.135$ velocities in the source Table~2 are averages over the
tower-diameter group rather than measurements for this centre-panel condition,
so they are used only as contextual scales.  The source Fig.~11 is likewise
not imported: although it uses the same tower diameter, its initial head and
water level are $0.305$ and $0.254$~m, respectively, and therefore define a
different test.

The quantitative one-dimensional history uses a crown-datum pressure and a
pressure-driven tracked-front application closure. Its
Test~B mouth resistance and coverage
parameters were selected from a predeclared six-point bounded screen:
$K_{\mathrm{slug}}/K_{\mathrm{g}}=1.25$, 1.50, or 1.75, and
$\alpha_{g,\mathrm{j}}^{\mathrm{full}}=0.10$ or 0.12. The admissible 1.50/0.10
pair minimized the native interface error. The resulting errors are in-sample
closure diagnostics rather than an independent calibration test. No time
shift, curve fit, or level-trajectory smoothing is introduced.
The spatial snapshots serve a different role from the quantitative history.
Figure~\ref{fig:caseB_snapshots} uses a separate visualization archive. Its
horizontal state is Tosan-based, and its tower is reconstructed from saved
$Y_{fs}$, $Y_{int}$, and jet-height scalars. It supports morphological
comparison rather than quantitative tower-field validation. The planar
OpenFOAM calculation also serves only as an independent spatial diagnostic. It
retains the published pipe lengths and initial heads, but uses the
area-equivalent slit width $W=D_t^2/D$ and neglects surface tension. It is a
controlled two-dimensional surrogate for the release pathway, not a
geometry-exact representation of the $12.7$-mm circular tower.

The quantitative one-dimensional archive selects the observed geysering
branch. The standard two-dimensional archive selects the same branch. In the
one-dimensional result, the shock-fitted gas front crosses
$Y^{*}_{int}=0.05$ at $T_{\mathrm{rel}}^{*}=3.622$. The conservative
free-surface observer reaches the predeclared one-cell rim-completion level at
$T_{\mathrm{rel}}^{*}=3.923$, where $1-\Delta z/L=0.984$. This occurs before the
computed front/free-surface catch at $T_{\mathrm{rel}}^{*}=4.051$. The
corresponding two-dimensional entry and rim events occur at
$T_{\mathrm{rel}}^{*}=3.85$ and $4.11$. Both calculations therefore place rim
arrival ahead of the resolved upper-interface trajectory, although the
two-dimensional surrogate develops the sequence later. The archived
two-dimensional window does not resolve a subsequent gas-core catch, and none
is inferred from this comparison. The three rows in
Fig.~\ref{fig:caseB_snapshots} compare the complete pipe--tower domain at the
same physical times without event alignment: gas-pocket approach at
$t=6.35$~s, tower entry and water-column lift at $7.25$~s, and the initial
discharge/venting response at $7.70$~s. This common clock exposes the timing
difference; the rows do not supply the event metrics quoted above.

\begin{figure}[htbp]
\centering
\includegraphics[width=\textwidth]{caseB_1d2d_snapshots_3frame_current.pdf}
\caption{Test~B phase evolution at $t=6.35$, $7.25$, and $7.70$~s:
present 1D model (left) and planar OpenFOAM simulation (right). Blue denotes
water and white denotes gas.}
\label{fig:caseB_snapshots}
\end{figure}

The native-time histories in Fig.~\ref{fig:caseB_curves} separate close
one-dimensional agreement in the bulk scales from larger event-timing and
two-dimensional biases. The experimental pre-eruption stand is
$Y^{*}_{fs}=0.823$, compared with $0.834$ in the one-dimensional model and
$0.869$ in the two-dimensional surrogate. At the crown datum, the
one-dimensional pressure plateau is $H^{*}=0.754$, versus the digitized
experimental pseudo-median value $0.758$ used in
Table~\ref{tab:caseB_summary}. The black curve itself is the pointwise
arithmetic mean of the three digitized repetitions. The standard raw
two-dimensional pre-release mean is higher, $H^{*}=0.838$ over
$1\le T_{\mathrm{rel}}^{*}\le3$. The displayed 0.40-s one-dimensional pressure
mean falls below $H^*=0.30$ at $T_{\mathrm{rel}}^{*}=4.097$, within the
experimental collapse envelope. The standard two-dimensional head crosses the
same threshold at $4.823$. All of these comparisons retain the native time
coordinate. For the shape-only comparison in
Fig.~\ref{fig:caseB_curves}(a), the blue pressure history comes from a separate
full-open-rim two-dimensional archive. It is displayed after one constant
offset, $\Delta H^{*}=-0.104$; no time shift or shape fit is applied. This
curve is excluded from the raw quantitative metrics. All quoted
two-dimensional pressure values and Table~\ref{tab:caseB_summary} use the
standard, unshifted-amplitude archive.

\begin{figure}[htbp]
\centering
\includegraphics[width=\textwidth]{caseB_experiment_1d2d_curves.pdf}
\caption{Test~B comparison with \cite{vasconcelos2011}: (a) normalized
pressure head and (b) normalized free-surface and interface elevations.
Black denotes the experiment, orange the present 1D model, and blue OpenFOAM.}
\label{fig:caseB_curves}
\label{fig:caseB_pressure}
\label{fig:caseB_levels}
\end{figure}

The level histories recover the rising curvature, but the tracked interface
remains too early and high. The isolated digitized point at
$(T_{\mathrm{rel}}^{*},Y^{*})=(3.847,0.92)$ is retained in the plot for
provenance. It is not used as an arrival estimate because it lies
outside all three rising-interface tracks. Zero-level baseline/sentinel symbols
outside the rising tracks are also excluded from the trajectory comparison and
are not treated as positive-height measurements. With these symbols excluded,
the native-time rising-track RMSEs are $0.042$ and $0.206$ for the
one-dimensional $Y^{*}_{fs}$ and $Y^{*}_{int}$ histories, respectively; their
signed biases are $+0.034$ and $+0.172$. Both computed tracks accelerate over
the rising bands, with late-to-early slope ratios of 1.37 and 1.96. The
corresponding two-dimensional RMSEs remain $0.052$ and $0.303$. The revised
native curvature therefore removes the former piecewise-linear shape error,
but not the early, high interface bias. The connected cellwise
$\alpha_g=0.50$ field front reaches only $0.016L$ in this calculation.
Accordingly, $Y^{*}_{int}$ denotes the computed shock-fitted front closure
rather than a resolved dense-core isocontour. No plotting or digitization
offset is used to remove the remaining separation.

\FloatBarrier

\begin{table}[htbp]
\centering
\caption{Native-time, unshifted-amplitude Test~B validation summary.
Experimental intervals are
digitized three-run spans from the centre panels of Figs.~6 and~8 in
\cite{vasconcelos2011}; scalar experimental values use the digitized
pseudo-median, whereas the black pressure curve in
Fig.~\ref{fig:caseB_curves}(a) is the pointwise three-run mean. The 2D column
uses the standard raw archive and excludes the display-offset blue pressure
curve. The one-dimensional pressure-collapse row uses the same 0.40-s
centered mean shown in panel (a); its level-event rows use native samples.}
\label{tab:caseB_summary}
\begin{tabular}{lccc}
\toprule
Observable & Experiment & 1D & 2D \\
\midrule
Outcome & geyser, 3/3 & geyser & geyser \\
Pre-release $H^{*}$ & 0.758 & 0.754 & 0.838 \\
Pre-entry $Y^{*}_{fs}$ & 0.823 & 0.834 & 0.869 \\
$Y^{*}_{int}=0.05$ crossing, $T_{\mathrm{rel}}^{*}$ & 3.648--3.742 & 3.622 & 3.848 \\
First rim arrival, $T_{\mathrm{rel}}^{*}$ & 3.900--3.981 & 3.923$^{\mathrm{a}}$ & 4.114 \\
Pressure pseudo-median $H^{*}<0.30$, $T_{\mathrm{rel}}^{*}$ & 4.048 & 4.097 & 4.823 \\
\bottomrule
\end{tabular}
\vspace{0.3em}

\parbox{0.98\linewidth}{\footnotesize $^{\mathrm{a}}$The one-dimensional
rim event uses the predeclared one-cell completion criterion
$Y^{*}_{fs}\ge 1-\Delta z/L=0.9836$; the topology latch and plotted
geometric observer cross it at the same native sample.}
\end{table}

Table~\ref{tab:caseB_summary} reports only quantities that can be compared
without differentiating sparse video markers. The scalar experimental
pressure values are read from the digitized pseudo-median; the reported event
intervals retain the visible three-run scatter. Both calculations reproduce
the outcome and event ordering. Relative to the experimental intervals, the
one-dimensional tracked front begins slightly early and remains too high, but
the grid-complete rim event and pressure collapse fall within the measured
spans. The two-dimensional surrogate enters slightly late, spills late, and
retains its excess head longest.

\FloatBarrier

Across the two tests, the archived branch evidence selects venting without
overflow in the wide tower and eruption in the narrow tower. The Test~A
connected-pocket sensitivity recovers the pressure and free-surface scales but
delays terminal depressurization; its interface event times remain within the
three-run spread. It is not conflated with the
separate spatial reconstruction. For Test~B, the quantitative one-dimensional
history closely reproduces the pre-eruption pressure and water-column scales,
while the shock-fitted interface remains early and high. The grid-complete
rim event and pressure collapse lie within the experimental spans.
The standard planar calculation supports the same entry-to-overflow morphology,
but overpredicts the pre-release head and water level and delays the principal
events. Quantitative comparison remains anchored to the experiment. The
two-dimensional results and display reconstructions provide supporting evidence
subject to the limitations stated in their captions. Section~\ref{sec:discussion}
relates the remaining timing pattern to the shaft-mouth and entrainment
closures.

\FloatBarrier

\subsection{Campaign 2: entrapped-air release from a horizontal pipe into a
vertical riser (Cong, Chan and Lee, 2017)}
\label{sec:tests:cong2017}

The apparatus of \cite{cong2017} consists of a horizontal acrylic pipe with
diameter $D=0.05$~m and length approximately $6$~m. Its upstream end is
connected to a constant-head tank, and a tee-junction supports a vertical riser
of height $1.8$~m and variable diameter ($D_r=16$--$46$~mm). An air pocket of
prescribed initial length $L_0$ is created in the horizontal pipe between ball
valves. Release under the upstream head $H_0$ drives the pocket along the pipe
crown to the riser, where it vents through the initially water-filled column.
Video and high-speed cameras recorded the air--water interface trajectories in
the pipe and riser. Transducers measured pressure near the pipe end and at the
riser base (Fig.~\ref{fig:cong_apparatus}).

\begin{figure}[htbp]
\centering
\includegraphics[width=\textwidth]{campaign2_apparatus_redrawn.pdf}
\caption{Campaign 2 Series-B apparatus, redrawn from Fig.~1(b) of
\cite{cong2017}.}
\label{fig:cong_apparatus}
\end{figure}

The authors summarized the outcomes of their parameter study with
a two-parameter geysering criterion,
\begin{equation}
\frac{D_r}{D}\le 0.62
\qquad\text{and}\qquad
V^{*}_{air}=\frac{V_{air}}{(\pi D_r^2/4)\,H_0}\ge 3.42,
\label{eq:cong_criterion}
\end{equation}
which requires both a sufficiently narrow riser and a pocket volume large
relative to the riser water column.

This campaign evaluates the model as a regime classifier over a systematic
sweep, rather than against individual time histories as in Campaign~1. The
first comparison treats the seven high-speed-camera runs of the experimental
Series~B individually (B-H1\,--\,B-H7: $H_0=0.66$~m, $L_0=0.61$~m, riser
diameters $D_r/D=0.32$--$0.92$). Computed classification, pocket arrival time,
and mean rise velocity are compared with the measurements. The second
comparison covers a broader synthetic sweep over the experimental parameter
ranges ($D_r=16$--$46$~mm, $L_0=0.61$--$1.8$~m,
$H_0=0.66$--$0.88$~m; 63 configurations). The resulting classifications are
compared with the experimental criterion in Eq.~(\ref{eq:cong_criterion}).

All campaign runs use the same fully synchronous junction coupling as
Campaign~1 and the identical frozen solver configuration (the Campaign~1
Test~B closure set). No parameter is changed between campaigns or between
runs. Two geometric features are specific to this apparatus. The upstream
boundary is the constant-head tank, and the pocket is created
\emph{downstream} of the riser tee. The released pocket therefore travels
upstream along the crown toward the riser. The tee position, 2.88~m from the
tank, is fixed from the experiment's kinematics. The measured arrival times
and front speeds in Table~2 of \cite{cong2017} give a valve-to-tee distance
$T_a\,U_f$ of $2.51$/$1.92$/$1.32$~m for the three pocket lengths
$L_0=0.61$/$1.2$/$1.8$~m, all placing the tee at the same station. This
position is held fixed across the sweep. An earlier campaign version used a
staged horizontal/vertical treatment whose classifier reduced to a static
volume--diameter criterion. Those results were circular as a validation of
Eq.~(\ref{eq:cong_criterion}) and are superseded by the fully coupled results
reported here.

\subsubsection{Series B: classification, arrival times, and rise velocities}
\label{sec:tests:cong2017:seriesB}

The model reproduces five of the seven measured Series-B classifications
(Table~\ref{tab:cong_bh} and Fig.~\ref{fig:cong_bh}). It predicts geysering for
B-H1 ($D_r/D=0.32$) and no geysering for B-H4--B-H7
($D_r/D=0.62$--$0.92$), with no false positives. The two intermediate
geysering runs, B-H2 and B-H3 ($D_r/D=0.42$, $0.52$), are misclassified as
non-geysering. Their computed free surfaces rise to $0.87$ and $0.82$~m but do
not reach the rim. The error is one-sided: the model under-predicts eruption.
Section~\ref{sec:discussion} examines how volume-neutral mouth exchange may
contribute to this bias. The same exchange allows the narrow-riser pocket to
rebuild the overpressure that powers the B-H1 geyser, but it also cancels part
of the one-to-one column lift in the reservoir-fed layout as riser width
increases. The computed arrival times are late by $1.5$~s on average, with the
largest delay in the widest riser where the arrested crown current covers the
final reach slowly. The measured spread is $7.84$--$8.22$~s.

The rise-speed comparison reveals an overly pulsed computed response. In the
geysering run B-H1, the computed eruption-window climb speeds
($v_{fs}=1.71$, $v_{int}=1.28$~m/s as the fastest sustained $0.6$-s climb)
exceed the respective measured whole-event means ($0.92$, $1.23$~m/s). The
interface speed is close, but the free-surface speed is markedly high. The
eruption jet is therefore over-sharp, consistent with the narrow-riser bias
found in Test~B of Campaign~1. In the no-geysering branch, the computed gas-nose climb speeds
($0.69$--$1.20$~m/s) exceed the measured $0.42$--$0.48$~m/s. The computed
climb is concentrated in surge pulses rather than the observed steady
Taylor-scale ascent. The computed free-surface response
($0.20$--$0.65$~m/s during the pulses, zero sustained rise) has the same pulsed
character, compared with the measured steady $0.2$--$0.26$~m/s.

The synchronous coupling recovers the measured ranking of the pipe-end
pressure transients. The geysering run B-H1 develops a cycle-averaged
compression surge to $1.71\,H_0$, compared with the reported value of about
$1.9\,H_0$. The no-geysering run B-H6 shows only a gentle post-arrival hump to
$1.02\,H_0$, compared with the reported value of about $1.4\,H_0$. The staged
treatment gave both runs the same $2.07\,H_0$ peak. The synchronous coupling
instead recovers the experimental differentiation: a wide riser vents the
pocket and thereby damps the pipe-end pressure. This gain is accompanied by
the two intermediate-diameter classification misses discussed above.

\begin{table}[htbp]
\centering
\caption{Campaign 2, experimental Series B of \cite{cong2017}
($H_0=0.66$~m, $L_0=0.61$~m): measured against computed classification,
pocket arrival time $T_a$, maximum computed free-surface height
$Y_{fs,\max}$ (riser top at $1.8$~m), and rise speeds of the free surface
($v_{fs}$) and gas nose ($v_{int}$). Measured speeds are whole-event means
(Table~2 of \cite{cong2017}); computed speeds are the fastest sustained
$0.6$-s climb in the venting window. Misclassified runs are marked
$^{\dagger}$.}
\label{tab:cong_bh}
\resizebox{\textwidth}{!}{%
\begin{tabular}{lcccccccccc}
\toprule
& & & \multicolumn{2}{c}{Geysering} & \multicolumn{2}{c}{$T_a$ (s)}
& & \multicolumn{2}{c}{$v_{fs}$ (m/s)} & $v_{int}$ (m/s) \\
\cmidrule(lr){4-5}\cmidrule(lr){6-7}\cmidrule(lr){9-10}
Run & $D_r/D$ & $V^{*}_{air}$ & Meas. & Model & Meas. & Model
& $Y_{fs,\max}$ (m) & Meas. & Model & Meas. / Model \\
\midrule
B-H1 & 0.32 & 9.03 & yes & yes & 8.07 & 7.56 & 1.80 & 0.92 & 1.71 & 1.23 / 1.28 \\
B-H2$^{\dagger}$ & 0.42 & 5.24 & yes & no & 7.84 & --   & 0.87 & 0.77 & --   & 1.02 / -- \\
B-H3$^{\dagger}$ & 0.52 & 3.42 & yes & no & 8.18 & 9.56 & 0.82 & 0.66 & 0.43 & 0.92 / 1.24 \\
B-H4 & 0.62 & 2.40 & no & no & 8.14 & 9.40  & 0.82 & 0.21 & 0.53 & 0.42 / 1.20 \\
B-H5 & 0.72 & 1.78 & no & no & 8.10 & 9.88  & 0.86 & 0.26 & 0.65 & 0.48 / 1.05 \\
B-H6 & 0.82 & 1.37 & no & no & 8.10 & 9.16  & 0.78 & 0.25 & 0.39 & 0.48 / 1.05 \\
B-H7 & 0.92 & 1.09 & no & no & 8.22 & 11.44 & 0.79 & 0.20 & 0.20 & 0.44 / 0.69 \\
\bottomrule
\end{tabular}
}
\end{table}

\begin{figure}[htbp]
\centering
\includegraphics[width=\textwidth]{cong2017_bh_series.png}
\caption{Campaign 2 Series B \cite{cong2017}: (a) geysering classification,
(b) pocket-arrival time, (c) free-surface rise speed, and (d) gas-nose rise
speed versus $D_r/D$.}
\label{fig:cong_bh}
\end{figure}

The two signature runs clarify how riser diameter changes the pressure and
water-column response (Fig.~\ref{fig:cong_signature}). These runs are also
examined as individual histories in
Sections~\ref{sec:tests:cong2017:seriesB} and the case folders. Their archived
reproductions include per-run comparisons with the digitized level and pressure
histories from Figs.~7, 9 and~10 of \cite{cong2017}. In B-H1
($D_r/D=0.32$), the released pocket arrives at $7.56$~s (measured $8.07$~s).
The entering gas drives the free surface upward in two surge pulses, while
arrest of the supply-line counter-flow compresses the pocket to
$1.71\,H_0$ (cycle-averaged; reported $\approx1.9\,H_0$). The resulting column
ejection places the free surface at the rim while the gas nose remains below
it, which is the geysering signature. In B-H6 ($D_r/D=0.82$), the same release
delivers the same pocket, but the wide riser passes the gas with only a gentle
free-surface rise to $0.78$~m. The pocket vents without eruption, and its
pressure never exceeds $1.02\,H_0$ after arrival. Because only $D_r$ changes,
the pipe-end pressure provides the experiment's cleanest controlled contrast.
The model reproduces its sign and ranking, although both pressure levels are
under-predicted. The $1.4\,H_0$ no-geyser hump is missed because the computed
free-surface lift is under-predicted in the reservoir-fed layout
(Section~\ref{sec:discussion}).

\begin{figure}[htbp]
\centering
\includegraphics[width=\textwidth]{cong2017_signature.png}
\caption{Campaign 2 signature cases \cite{cong2017}: free-surface and gas-nose
trajectories for (a) B-H1 and (b) B-H6, and (c) their pocket-pressure heads.}
\label{fig:cong_signature}
\end{figure}

The common-clock spatial comparison for B-H1 supports the same branch but
exposes different eruption clocks (Fig.~\ref{fig:cong_bh1_1d2d}). At the
physical times $8.5$, $12.8$, and $13.0$~s, the frozen 1D archive reports a
much earlier free-surface rise. It reaches $1.79$~m above the pipe crown by
$13.0$~s, whereas the refined planar OpenFOAM field still places the water
column below the $1.8$~m rim. The 2D calculation reaches $98\%$ of the rim at
$13.865$~s and resolves water above $2.1$~m at $14.8$~s. Both descriptions
select the measured geysering branch, but their eruption clocks and spatial
representations differ. The 2D gas-arrival time, $8.50$~s, is close to the
measured $8.07$~s. Its subsequent eruption is nevertheless late relative to
the reported $9.55$~s, and its post-arrival pressure surge is not reproduced.
The sequence therefore provides supporting evidence for branch selection and
morphology, not a quantitative OpenFOAM validation.

\begin{figure}[htbp]
\centering
\includegraphics[width=\textwidth]{cong2017_bh1_1d2d_phase_sequence.png}
\caption{B-H1 riser phase sequence at $t=8.5$, $12.8$, $13.0$, and
$14.8$~s: present 1D model (top) and planar OpenFOAM simulation (bottom).
The archived 1D run ends at $13.0$~s.}
\label{fig:cong_bh1_1d2d}
\end{figure}

The complete-domain comparison for B-H6 preserves the observed non-geysering
sequence at the photographic clock (Fig.~\ref{fig:cong_bh6_multiframe}). The
three times correspond to gas arrival ($8.7$~s), developed rise ($9.9$~s), and
late venting ($10.9$~s). Each row places the published B-H6 riser photograph
between the full 1D and planar OpenFOAM domains, with no event alignment. The
1D reconstruction passes a compact gas pocket with weaker water-column lift,
whereas the 2D VOF field gives a stronger and more spatially distributed late
response. Because the high-speed record images the riser rather than the full
horizontal pipe, the photographic column validates the riser sequence. The
outer panels show the corresponding system-level phase distribution.

\begin{figure}[htbp]
\centering
\includegraphics[width=0.98\textwidth]{cong2017_bh6_experiment_1d2d_full_domain_3frame.pdf}
\caption{B-H6 full-domain comparison at (a) $8.7$, (b) $9.9$, and
(c) $10.9$~s: present 1D model (left), experiment \cite{cong2017} (centre),
and planar OpenFOAM simulation (right).}
\label{fig:cong_bh6_multiframe}
\end{figure}

The aggregate Series-B comparison conceals how the two descriptions differ
within an individual non-geysering event. Figure~\ref{fig:cong_bh6_1d2d}
therefore overlays the digitized B-H6 markers from Fig.~7(a) of
\cite{cong2017} with a geometry-matched one-dimensional application and the
area-equivalent planar OpenFOAM calculation. All traces retain physical time;
only the dense two-dimensional trace receives a centred $0.05$-s running
median for display. Both calculations reproduce the observed non-geysering
outcome. The two-dimensional gas reaches the riser at $8.04$~s, close to the
measured $8.10$~s, whereas the one-dimensional arrival is $8.66$~s. The
trajectory comparison is less complete: free-surface marker RMSE is
$0.097$~m for the two-dimensional result and $0.310$~m for the
one-dimensional result, while gas-nose RMSE is $0.245$ and $0.317$~m,
respectively. Both under-predict the final lift, but the spatial calculation
substantially reduces the free-surface error.

The focused one-dimensional trace is a geometry- and valve-matched application
variant whose horizontal gas front uses the Benjamin-celerity closure; it is
shown to diagnose this case and does not replace the frozen Series-B values in
Table~\ref{tab:cong_bh}. Likewise, the planar calculation is supporting
evidence for spatial consistency rather than an expansion of the primary
validation scope.

\begin{figure}[htbp]
\centering
\includegraphics[width=\textwidth]{cong2017_fig7a_bh6_experiment_1d2d.pdf}
\caption{B-H6 free-surface ($Y_{fs}$, solid) and interface ($Y_{int}$,
dashed) trajectories overlaid on Fig.~7(a) of \cite{cong2017}. Dark red and
green denote the present 1D model and OpenFOAM, respectively.}
\label{fig:cong_bh6_1d2d}
\end{figure}

A source-matched comparison preserves the experimental parameter context by
transcribing the B-H6 row of Table~2 in \cite{cong2017}
(Table~\ref{tab:cong_bh6_source_matched}). The second block reports only
quantities available robustly from all three records. This makes the outcome
and arrival-time comparison explicit without substituting peak-based
simulation values for the source table's whole-event velocity definitions.

\begin{table}[htbp]
\centering
\caption{Source-matched B-H6 parameter and result comparison. The first block
transcribes the B-H6 row of Table~2 in \cite{cong2017}; the second block adds
the experimental, 1D, and OpenFOAM 2D outcome, gas-arrival time, peak levels,
and marker RMSE. Dashes denote quantities not defined for the experimental
reference.}
\label{tab:cong_bh6_source_matched}
\small
\begin{tabular}{@{}lclc@{}}
\toprule
\multicolumn{4}{c}{Published B-H6 condition and whole-event measures} \\
\midrule
Parameter & Value & Parameter & Value \\
\midrule
$D_r$ (m) & 0.041 & $H_0$ (m) & 0.66 \\
$L_0$ (m) & 0.61 & $T_a$ (s) & 8.10 \\
$U_f/\sqrt{gD}$ & 0.443 & $D_r/D$ & 0.82 \\
$V_{\mathrm{air}}/V_w$ & 1.37 & Geyser & No \\
$\overline{v}_{fs}$ (m/s) & 0.246 & $\overline{v}_{int}$ (m/s) & 0.476 \\
$\overline{v}_{net}$ (m/s) & 0.235 & $v_{\mathrm{Taylor}}$ (m/s) & 0.219 \\
\bottomrule
\end{tabular}

\vspace{0.7em}

\begingroup
\footnotesize
\setlength{\tabcolsep}{2.5pt}
\begin{tabular}{@{}lcccccc@{}}
\toprule
Source & Outcome & $T_a$ (s) &
\shortstack{$Y_{fs,\max}$\\(m)} &
\shortstack{$Y_{int,\max}$\\(m)} &
\shortstack{RMSE $Y_{fs}$\\(m)} &
\shortstack{RMSE $Y_{int}$\\(m)} \\
\midrule
B-H6 experiment & No & 8.10 & 1.201 & 1.178 & -- & -- \\
1D & No & 8.66 & 0.786 & 0.700 & 0.310 & 0.317 \\
OpenFOAM 2D & No & 8.04 & 1.021 & 0.940 & 0.097 & 0.245 \\
\bottomrule
\end{tabular}
\endgroup
\end{table}

\FloatBarrier

\subsubsection{Regime map against the experimental criterion}
\label{sec:tests:cong2017:map}

The 63-configuration sweep yields a one-sided missed-eruption bias in the
criterion plane $(D_r/D,\ V^{*}_{air})$
(Fig.~\ref{fig:cong_map}). Each point represents one full synchronous
simulation classified without knowledge of Eq.~(\ref{eq:cong_criterion}); the
criterion lines are drawn only for comparison. The model agrees with the
experimental criterion in 39 of 63 configurations. In all 24 disagreements,
the criterion predicts geysering and the model does not. There are no false
positives. Every modelled eruption lies inside the criterion's geysering
region. The no-geysering half-plane $D_r/D>0.62$ and the small-volume region
$V^{*}_{air}<3.42$ are classified correctly throughout. The misses concentrate
at intermediate diameters ($D_r/D=0.42$--$0.62$) combined with large pocket
volumes ($L_0\ge1.2$~m). In these cases, the computed free surface rises to
$1.0$--$1.6$~m but remains below the $1.8$-m rim; seven of the 24 reach within
$0.45$~m of it. This pattern extends the Series-B under-lift bias across the
parameter plane. The model retains both end members and the interaction that
eruption requires narrow risers \emph{and} large relative pocket volume, but
requires more restrictive conditions for eruption and under-predicts it in the
transition band. For screening, the bias direction is consequential: the
frozen model misses geysers near the boundary rather than producing false
alarms.

\begin{figure}[htbp]
\centering
\includegraphics[width=0.85\textwidth]{cong2017_criterion_map.png}
\caption{Campaign 2 classification of 63 configurations in the criterion
plane of Eq.~(\ref{eq:cong_criterion}). Filled triangles denote model
geysering, circles non-geysering, and black rings disagreement with the
experimental criterion.}
\label{fig:cong_map}
\end{figure}

\FloatBarrier

\subsection{Campaign 3: stormwater geyser in a vertical shaft above a
junction chamber (Liu, Shao and Zhu, 2020)}
\label{sec:tests:liu2020}

The apparatus of \cite{liu2020} is a 1:20-scale model of a 27~m deep
Edmonton prototype dropshaft. A 5.80~m upstream pipe ($D=0.20$~m,
slope 1:100) feeds a transparent junction chamber
($0.30\times0.30\times0.45$~m) with an invert drop of $0.18$~m to a
5.95~m downstream pipe ($D_d=0.28$~m). A vertical shaft of diameter
$d_r=0.06$~m and height $1.22$~m is mounted at the chamber lid centre
and vents to the atmosphere. Inflow is imposed as a flow step
$Q_0\to Q_1$ with valve opening time $T_v\approx0.4$~s. Pressure
transducers PT1--PT4 monitor the shaft, the chamber lid (PT2, the
primary datum), the chamber floor, and the upstream crown
(Fig.~\ref{fig:liu_apparatus}).

\begin{figure}[htbp]
\centering
\includegraphics[width=\textwidth]{campaign3_apparatus_redrawn.pdf}
\caption{Campaign 3 apparatus and initial states, redrawn from Fig.~2 of
\cite{liu2020}.}
\label{fig:liu_apparatus}
\end{figure}

The three configurations use the documented baseline numerical settings from
the primary Campaigns~1--2 quantitative archives. The horizontal reach is a
single composite one-dimensional domain comprising the upstream pipe, chamber,
and downstream pipe. Internal momentum faces represent the chamber
connections, and a rigid-column ordinary differential equation for the shaft
water column is coupled at the lid tap. This reduced formulation does not
resolve three-dimensional chamber mixing, slug venting at the chamber gas port,
or the late-time arrival of a compressing air pocket along the upstream crown
at the level of the companion two-fluid/RTRP framework. These limitations are
stated explicitly where they affect the comparison.

\subsubsection{Branch pair A2/B3: downstream boundary flips the outcome}
\label{sec:tests:liu2020:AB}

Cases A2 and B3 share the same inflow step ($Q_0=20$, $Q_1=100$~L/s) and
differ only in the downstream tail: A2 discharges to an open channel
(weir-controlled, $h_d/D_d=1/4$), whereas B3 is initially surcharged
with a submerged tail gate. In the experiment A2 never geysered while
B3 produced a single-shoot eruption on every repetition
(Fig.~5 of \cite{liu2020}).

The unchanged model reproduces the A2/B3 branch flip. For A2, the riser column
remains below the rim ($h_{r,\max}=0.47$~m against a shaft height of
$1.22$~m). The steady chamber pressures agree with the digitized Fig.~3
record: PT3 is $5.26$~kPa against $4.99$~kPa (the model is about $5\%$ high),
and PT2 is $1.80$~kPa against $2.15$~kPa (the model is about $16\%$ low in the
standing-column head that
emerges from the composite domain rather than being prescribed). The computed
B3 run geysers, with the free surface reaching the shaft top and overflowing.
Its eruption timing is close, with peak PT2 at $1.68$~s against $1.47$~s.
However, the frozen elastic branch celerity $a_{wh}=40$~m/s underpredicts the
compressive peak amplitude: the PT2 peak is $31$~kPa against $55$~kPa
($-44\%$). A sensitivity run at $a_{wh}=80$~m/s gives $63.5$~kPa, showing that
the peak is highly sensitive to acoustic wave speed. This single sensitivity
test does not exclude a contribution from closure error. At the baseline
computed peak pressure, the experimental
regression in Fig.~7(a) gives approximately $2.50$~m
($h=0.694\,P_{\max}/\rho g+0.309$~m, $R^2=0.97$). The computed jet height is
$2.17$~m, about $0.33$~m lower.

\begin{figure}[htbp]
\centering
\includegraphics[width=\textwidth]{liu2020_a2_pressure.png}
\caption{Case A2 ($Q$: 20$\to$100~L/s, downstream open channel):
chamber and shaft pressures against time. Blue: digitized experimental
PT2/PT3 from Fig.~3 of \cite{liu2020}; red: present model.}
\label{fig:liu_a2_pressure}
\end{figure}

\begin{figure}[htbp]
\centering
\includegraphics[width=\textwidth]{liu2020_b3_pressure.png}
\caption{Case B3 ($Q$: 20$\to$100~L/s, downstream surcharged): PT1--PT3
pressure histories. Blue: experiment \cite{liu2020}; red: present model.}
\label{fig:liu_b3_pressure}
\end{figure}

\begin{figure}[htbp]
\centering
\includegraphics[width=0.72\textwidth]{liu2020_b3_h_pmax.png}
\caption{Case B3 on the eruption-height--peak-pressure diagram of
Fig.~7(a) of \cite{liu2020}. Open symbols: experimental family; filled
red circle: present model ($h=2.17$~m, $P_{\max}=31$~kPa at
$a_{wh}=40$~m/s).}
\label{fig:liu_b3_h_pmax}
\end{figure}

\subsubsection{Case C9: transient-driven geysers and the analytic
mass-oscillation model}
\label{sec:tests:liu2020:C9}

Case C9 ($Q_0=25\to Q_1=40$~L/s, surcharged downstream tail,
initial shaft column $h_{r0}=0.30$~m, sealed air wedge on the upstream
crown in the experiment) is the authors' detailed two-stage case: the
first two geysers are driven by the hydraulic transient before the
trapped pocket reaches the chamber at $t\approx6.46$~s; geysers
3--8 are pocket-driven and lie outside the present reduced solver.

For the phase-one comparison the pocket is omitted in the computation
(the all-water variant assumed in Eq.~(7) of \cite{liu2020}), while the
sealed-pocket run is retained as a qualitative cushioning check. The
analytic head at the chamber lid is
\begin{equation}
H_s(t)=h_{r0}+\frac{L_d}{gA_d}\frac{\mathrm{d}Q_u}{\mathrm{d}t}
\left(1-\cos\frac{2\pi t}{T}\right),
\qquad
T=2\pi\sqrt{\frac{L_d A_r}{gA_d}+\frac{h_{r0}}{g}},
\label{eq:liu_eq7}
\end{equation}
with $T=1.52$~s for C9.

The all-water model captures the first C9 pressure peak in both timing and
magnitude (Fig.~\ref{fig:liu_c9_phase1}). The figure compares the digitized PT2
record from Fig.~9, converted to head at the chamber-top datum, with
Eq.~(\ref{eq:liu_eq7}) and the all-water model over the first 5~s. The measured
first peak is $P_{1m}=10.69$~kPa at $t=0.50$~s, and the model gives
$11.5$~kPa at $0.54$~s. The free surface reaches the shaft top at $0.73$~s in
the experiment and $0.93$~s in the model. The oscillation period is
$T=1.45$~s in the paper (Eq.~(8)), compared with $1.13$~s in the model and
$1.52$~s in Eq.~(\ref{eq:liu_eq7}). The model period is closer to the analytic
value than to the measured period, but remains fast, consistent with the stiff
one-dimensional tail-gate coupling. After the first transient, the steady
heads agree within $\approx20\%$: the experimental PT2 final value is
$8.79$~kPa, and the model gives $10.6$~kPa. The
sealed-pocket variant qualitatively delays and cushions the first peak as in
the experiment. It then drifts after $\approx4$~s as the wedge compresses
through the reduced gas model, a phase-two process outside the stated scope.

\begin{figure}[htbp]
\centering
\includegraphics[width=\textwidth]{liu2020_c9_phase1.png}
\caption{Case C9, phase~1 ($t<5$~s): PT2 head $H_s$. Blue: experiment
\cite{liu2020}; orange dashed: Eq.~(\ref{eq:liu_eq7}); red: all-water model;
purple dotted: sealed-pocket variant.}
\label{fig:liu_c9_phase1}
\end{figure}

\begin{figure}[htbp]
\centering
\includegraphics[width=\textwidth]{liu2020_c9_riser.png}
\caption{Case C9: shaft water-column height $h_r$ for the all-water
(solid) and sealed-pocket (dotted) variants. The dashed line marks the
experimental time at which the free surface first reaches the shaft
top in phase~1.}
\label{fig:liu_c9_riser}
\end{figure}

\begin{table}[htbp]
\centering
\caption{Campaign~3 summary metrics. Acoustic peaks use the frozen
$a_{wh}=40$~m/s; no coefficient is retuned between cases.}
\label{tab:liu_summary}
\begin{tabular}{lccc}
\toprule
Quantity & A2 (no geyser) & B3 (geyser) & C9 phase~1 \\
\midrule
Geyser? & no / no & yes / yes & yes / yes \\
PT2 steady or peak [kPa] & $1.80$ / $2.15$ & $31$ / $55$ & $11.5$ / $10.7$ \\
Peak or top-arrival time [s] & -- & $1.68$ / $1.47$ & $0.54$ / $0.50$ \\
Shaft top reached [s] & -- & $\approx1.65$ & $0.93$ / $0.73$ \\
Oscillation period [s] & $0.30$ (PT2--PT3) & -- & $1.13$ / $1.45$ \\
\bottomrule
\multicolumn{4}{l}{\footnotesize Slash-separated cells: model / experiment; single values are identified by row.}\\
\end{tabular}
\end{table}

Across the three cases, the composite-domain formulation extends the frozen
Campaign~1--2 solver to a prototype-motivated junction geometry. The A2/B3
pair reproduces the change in regime caused by the downstream boundary, and C9
phase~1 reproduces the transient mass-oscillation geyser. Agreement covers
regime selection and timing, as well as amplitude where the governing response
is hydraulic rather than acoustic. The main residuals are the
under-resolution of short compressive pressure spikes with the frozen elastic
celerity and the deliberate omission of pocket-transport physics after the
pocket reaches the chamber.
